%% ===========================================================================
%%  REV
%%  Rethinking the Execution Venue for RDMA Disaggregated-Memory Indexes
%%
%%  Compile twice:  pdflatex paper && pdflatex paper
%%
%%  STRUCTURE
%%    1  Introduction
%%    2  Background
%%    3  Two levers, and how we measure them
%%    4  DEX
%%    5  CHIME
%%    6  Discussion
%%    7  Related work
%%    8  Conclusion
%%
%%  PROVENANCE. All measured, no projections, no cross system numbers.
%%    dex/build/results/summary.csv, summary_full.csv, and the per run
%%        [ALL OPS] blocks in dex_*_offload-{off,on}_cache*mb.log
%%        totalThreadCount 36, memThreadCount 4, 50 M keys, scan length 100
%%    dex/build/results/*_mt{2,4,6}.log     memory thread sweep
%%    CHIME/results/leafstudy2_compute.csv  lever sweep, 34 threads per node,
%%        16 B values, 50 M keys, 30 M ops, scan length 100
%%    CHIME/results/RANGE_SCANS.md          leaf cache first cut, pre repair
%%  Figures are emitted by paper/make_figures.py into paper/figures/.
%%
%%  House style: no dash characters in the prose. No code identifiers in the
%%  body text; mechanisms are described in design terms.
%%
%%  TODO before submission: bibliography. Systems are named in prose and the
%%  \cite keys are not yet wired.
%% ===========================================================================
\documentclass[sigconf,nonacm]{acmart}

\usepackage{graphicx}
\usepackage{subcaption}
\usepackage{booktabs}
\usepackage{amsmath}
\usepackage{xcolor}
\usepackage{balance}

\graphicspath{{./figures/}{../hybrid_plots/}%
  {C:/Users/abhir/Downloads/REV/paper/figures/}}

\settopmatter{printacmref=false}
\renewcommand\footnotetextcopyrightpermission[1]{}
\pagestyle{plain}
\emergencystretch=2em

\begin{document}

\title{Rethinking the Execution Venue for Indexes\\on RDMA Disaggregated Memory}
\subtitle{A cache changes how often you pay the toll, never what the toll is}

\author{}
\affiliation{\institution{}\country{}}

% ===========================================================================
\begin{abstract}
Disaggregated memory splits a cluster into CPU rich compute nodes and capacity
rich, CPU poor memory nodes joined by RDMA. Every ordered index on that substrate
must answer a question that does not arise in a shared memory system: \emph{where
does an index operation execute}. The answer is almost always the same, and it is
almost always fixed at design time. The index caches part of itself on the
compute node and leaves the memory node idle, because memory side CPU is assumed
too scarce to be a participant.

We show by measurement, on two state of the art disaggregated B$^+$ trees, that
this leaves a structural gap. A cache changes how \emph{often} an operation
crosses the network and never what a crossing \emph{costs}, so performance is
governed by whether the working set fits, and range scans, which stream through
leaves they never revisit, cannot be helped by cache budget at all. In DEX an
eightfold cache increase buys uniform scans a factor of 1.32 in throughput, and
the rate at which its leaf path misses does not fall across that sweep by any
amount we can measure. In CHIME, whose compute cache holds internal nodes only,
throughput falls by 3.7 times in a single step when the internal index stops
fitting.

We then make the execution venue a runtime decision, gated so that the memory
node is engaged only after the compute side cache has already failed, and we
measure each system against itself. Thirteen of sixteen DEX workload cells then
exceed, simultaneously on throughput and tail, the best that design reaches
anywhere in an eightfold cache sweep. On scans the gated configuration at one
eighth of the memory beats the cached configuration at full memory on both axes
at once. In CHIME the two levers separate cleanly: compute side caching raises
the plateau where the index fits and does almost nothing below it, while memory
side execution removes the cliff where it does not and is exactly inert above it.
Used together they exceed what either reaches alone. Finally we report a
deliberate falsification: an optimisation that shrank the bytes per covered leaf
roughly fortyfold while leaving the number of crossings unchanged lost 41\% of
throughput, and lost 15\% even at a 59\% cache hit rate. Given a path that
removes crossings instead, the same mechanism gains 19\%.
\end{abstract}

\maketitle

% ===========================================================================
\section{Introduction}
\label{sec:intro}

An index on disaggregated memory has two levers against remote cost, and the
literature has largely treated them as rivals.

A \textbf{compute side cache} decides \emph{how often} an operation touches the
network. It is attractive because it is portable: it assumes nothing of the
memory node beyond a region of registered memory, so it survives the strictest
reading of disaggregation, in which the memory node has almost no usable CPU at
all.

\textbf{Memory side execution} decides \emph{what a touch costs}. It runs index
logic on memory node CPUs so that a request returns a result rather than bytes.
It is the lever the literature reaches for last, and the reason is the same
axiom: the memory node is assumed too CPU poor to participate.

Two recent designs illustrate how much can be extracted from the first lever
alone, and they do it in different directions. DEX reduces how often a remote
trip happens, with a path aware compute side cache, pointer swizzling, a scalable
replacement structure and a probabilistic admission policy. CHIME reduces how
much a trip carries and how many trips a search needs, pairing a hopscotch leaf
layout with a compute side cache so that one one sided read covers more of the
search. Both are careful, and both are complete answers within their assumption.

This paper asks what that assumption costs, and it asks it by measurement.

\paragraph{The finding.}
The quantity these designs spend their budget on is not the quantity that bounds
them. A cache can make the expensive path rarer. It cannot make it cheaper, and
on this fabric the cost of a crossing is nearly independent of how many bytes it
carries: a seventy byte read and a five hundred byte read cost about the same.
That has three consequences we measure directly.

First, the caching lever \emph{saturates, and reports its own saturation}. When a
leaf level remote procedure path is available, the rate at which it fires counts
exactly how often the cache failed to resolve the leaf. In DEX under uniform
access that rate is flat: every lookup misses its leaf at the first three cache
points and 99.7\% still do at the fourth, and uniform scans hold at 1.3748 calls
per operation to four decimal places across the entire eightfold sweep. More
memory is not producing a diminishing return there. It is producing no return.

Second, where the cache does have a residency boundary to cross, the design falls
off it. CHIME caches internal nodes only, and its internal index is roughly 90 to
100 megabytes at the geometry we use. Above that budget its throughput is flat and
good; one step below it, uniform point throughput falls by a factor of 3.7.

Third, both effects are removed by the same change, and it is a small one:
convert a leaf level cache miss into a request that the memory node answers, and
gate that conversion on the miss so that the memory node is engaged only when the
cache has already failed. In DEX this puts thirteen of sixteen workload cells past
the best that design reaches anywhere in an eightfold cache sweep, on throughput
and tail simultaneously. In CHIME it removes the cliff entirely while remaining
exactly inert above the boundary.

\paragraph{Levers, not knobs.}
The two levers turn out to be complementary in a regime sense rather than
substitutes. Caching raises the plateau where the working set fits and does
almost nothing below it. Memory side execution removes the cliff below the
boundary and contributes nothing above it. They act on disjoint parts of the same
curve, which is why a system that owns both and never asks which regime it is in
can still collapse, and why the pair reaches an operating point that neither
reaches alone.

\paragraph{Contributions.}
\begin{itemize}
  \item A measurement study of two state of the art disaggregated B$^+$ trees
    that locates, in each, the point at which the caching lever stops returning
    value, and shows that each design already carries a signal reporting it
    (Sections~\ref{sec:dex} and~\ref{sec:chime}).
  \item \emph{Run at a time scan pushdown}: a range scan served in a constant
    number of network round trips rather than one per covered leaf, by executing
    the leaf chain traversal at the memory node and returning a run of results in
    a single reply. Scans are the case where the caching lever is structurally
    unavailable, and no prior disaggregated index offloads one as a run
    (Section~\ref{sec:dex:add}).
  \item \emph{Encoding agnostic memory side execution}: the demonstration that
    pushdown requires only that the memory node speak the index's own node
    encoding, not that the index store plain records. We establish this by
    instantiating one execution path on two incompatible encodings, one of which
    exists precisely to make one sided reads safe without memory side CPU
    (Section~\ref{sec:chime:add}).
  \item \emph{The lever map}: a measurement separating the two levers by regime,
    showing that one owns the plateau and the other the cliff, and that their
    combination exceeds the unmodified design's own ceiling on a quarter of the
    memory it needs to reach that ceiling (Section~\ref{sec:chime:map}).
  \item A falsification test of the underlying account. Crossings and not bytes
    are the currency, which forbids a specific and tempting optimisation. We built
    that optimisation, measured the predicted loss, explained it, and repaired it
    (Section~\ref{sec:chime:falsify}).
\end{itemize}

% ===========================================================================
\section{Background}
\label{sec:bg}

\paragraph{Disaggregated memory.}
A disaggregated cluster separates compute nodes, which have many cores and little
memory, from memory nodes, which have much memory and few cores, and joins them
with RDMA. One sided verbs let a compute node read or write remote memory without
involving the remote CPU; two sided verbs require a thread on the far side to
receive and act. Because memory side cores are the scarce resource, one sided
access is the default, and a design that requires memory side threads has to
justify them.

\paragraph{Why disaggregated indexes cache inner nodes.}
An index operation walks from a root to a leaf, and each level not held locally
costs a round trip. Caching the upper levels is unusually attractive on this
substrate because inner nodes need no coherence: a stale inner node sends the
search to the wrong leaf, which the operation detects and retries, so a compute
node may hold an out of date copy without returning a wrong answer. Leaves have
no such property, and several designs therefore exclude them from the cache
rather than pay for cross node invalidation.

\paragraph{The two systems.}
\emph{DEX} caches both inner and leaf pages under a path aware policy, so that a
cached path is contiguous from the root downward, and admits leaves only
probabilistically because leaf residency scales with the data rather than with
the index. It offloads point operations at the cache boundary under a runtime cost
model, and it decomposes a range scan spanning several leaves into repeated single
key lookups, because it deliberately maintains no leaf links.

\emph{CHIME} caches internal nodes only and spends its remaining design budget on
making the one unavoidable remote leaf access smaller: a hopscotch leaf layout
confines a key to a short neighbourhood so a reader fetches a segment rather than
a whole leaf, a hotspot structure remembers where a hot key sits so the read
shrinks to a single entry, and a replicated per group metadata copy makes a
partial read self describing. Its leaves also carry per cacheline version bytes,
so a one sided reader can detect a torn read without taking a lock.

Every one of CHIME's mechanisms reduces bytes. None reduces crossings. That
distinction is the subject of this paper.

% ===========================================================================
\section{Two Levers, and How We Measure Them}
\label{sec:method}

\subsection{The levers}

Table~\ref{tab:levers} states the two levers and, more importantly, where each
one fails. Neither lever is new. What we report is that their failure modes are
disjoint, so that the choice between them is a question about regime rather than
a question about which is better.

\begin{table}[t]
\centering
\small
\caption{The two levers against remote cost. The row that matters is the third.}
\label{tab:levers}
\begin{tabular}{@{}p{0.13\linewidth}p{0.37\linewidth}p{0.37\linewidth}@{}}
\toprule
 & compute side caching & memory side execution \\
\midrule
removes & crossings for data that \emph{fits} & the \emph{cost} of a crossing
  that happens anyway \\
\addlinespace
scales with & cache budget against working set size & memory side CPU \\
\addlinespace
\textbf{fails when} & the working set does not fit, has no hot core, or is
  streamed rather than revisited & the cache is already succeeding, when it is
  pure overhead \\
\bottomrule
\end{tabular}
\end{table}

\subsection{Every comparison is a system against itself}

Both studies are ablations. The two configurations in every figure are one
runtime setting apart on one binary, running the same workload generator against
the same tree with the same thread count on the same hardware. Nothing else
differs, so every movement is causally attributable to the setting.

We deliberately report no cross system numbers. Such a comparison must equalise
the dataset, the cache accounting, the thread count and the latency definition
across codebases that measure each differently, and every one of those
equalisations is a place where an effect can be manufactured or hidden. A per
thread cache and a shared cache are not the same quantity; a system driving
clients from one machine and a system running the same binary on two nodes do not
mean the same thing by a thread count; a mean and a tail are not interchangeable.
The claims here are factors of two to six between two configurations of one
binary, and a comparison whose thread accounting is off by a factor of two can
move a ratio by as much as the effect being claimed. Section~\ref{sec:limits}
states what we would fix before adding one.

\subsection{What replaces an external reference}

Removing the external baseline removes a line from a chart, so something must say
whether a movement is large. Three quantities do, and all three are internal to
the system under test.

\begin{enumerate}
  \item \textbf{The design's own ceiling.} The best value the unmodified design
    reaches at \emph{any} cache size in the sweep. A configuration that exceeds
    that ceiling has done something more memory could not do.
  \item \textbf{Cache equivalence.} How much cache the unmodified design would
    need to reach what the modified one reaches at the smallest budget swept.
    This is the provisioning question stated without reference to anyone else.
  \item \textbf{Crossings per operation.} Network round trips, counting both one
    sided reads and remote procedure calls. This is what the fabric charges for,
    it is measured directly rather than inferred from throughput, and it is the
    mechanism behind the other two.
\end{enumerate}

\subsection{Setup}

Both studies run on the same two node RDMA cluster, one compute node and one
memory node. The workload matrix is four cells in each: point lookup and hundred
key range scan, crossed with uniform and Zipfian key distributions at skew 0.99.
Both bulk load fifty million keys and measure a read only workload, so nothing
here bounds insert, update or delete behaviour.

DEX runs thirty six compute threads against four memory side service threads and
sweeps its compute side cache across 64, 128, 256 and 512 megabytes, giving
sixteen cells per configuration. Latency is read from the per operation
histogram in each run's own log rather than from an aggregate, so a tail is a
measured tail.

CHIME runs thirty four application threads per node against four memory side
threads. Its sweep is larger because it has two levers to cross: four cache
budgets against leaf caching off and on against memory side execution off and on,
over the same four workloads, for sixty four cells and roughly six hours per
node. Its leaf cache is carved \emph{out} of the total budget rather than added to
it, so every arm occupies identical compute side memory and caching leaves must
earn its share against the inner nodes it displaces.

% ===========================================================================
\section{DEX}
\label{sec:dex}

\subsection{The design as built}

DEX pulls the caching lever hard and well. Its compute side cache holds both
inner and leaf pages under a path aware policy, so that the resident portion of
a root to leaf path is contiguous and a traversal degenerates into local pointer
chasing for as long as the path stays resident. Replacement avoids the
centralised structures that would serialise under the replacement rate a
disaggregated setting produces. Admission is probabilistic and asymmetric:
inner nodes are always admitted, leaves only one time in ten, because leaf
residency scales with the dataset rather than with the index and admitting a leaf
usually evicts something hotter.

DEX also has the second lever, and this is what makes it the more instructive
subject. It offloads point operations at the cache boundary when a runtime cost
model judges a remote procedure call cheaper than the remaining descent. What it
does not do is treat the two levers as covering different regimes. Both decisions
are evaluated locally, at a single node miss, and nothing in the design observes
whether the caching lever is succeeding for an operation class as a whole.

Range scans are the sharpest case. DEX maintains no leaf links, a simplification
that keeps its pointer management uniform, and a scan spanning several leaves is
therefore decomposed into repeated single key lookups. Every covered leaf pays its
own descent, its own cache probe and its own offload decision.

\subsection{The caching lever saturates, and says so}
\label{sec:dex:sat}

The first measurement is the cache sweep with memory side execution absent.
Across an eightfold increase, from 64 to 512 megabytes, uniform scans move from
0.42 to 0.56 million operations per second and uniform lookups from 1.29 to 2.52.
Tail latency moves less than throughput on every workload, and much less on
scans: the two scan workloads never get their 99th percentile below 92
microseconds anywhere in the sweep.

\begin{table}[t]
\centering
\small
\caption{What eight times the memory buys DEX. ``Ceiling'' is the best value
reached anywhere in the 64 to 512\,MB sweep. Tail improves less than throughput
on every workload, and on scans by very little.}
\label{tab:envelope}
\begin{tabular}{llrrr}
\toprule
 & & ceiling & throughput & tail \\
op & dist & Mops & gain & gain \\
\midrule
lookup & uniform & 2.52 & $1.96\times$ & $1.60\times$ \\
lookup & zipf    & 5.06 & $2.47\times$ & $1.66\times$ \\
scan   & uniform & 0.56 & $1.32\times$ & $1.21\times$ \\
scan   & zipf    & 0.95 & $1.59\times$ & $1.23\times$ \\
\bottomrule
\end{tabular}
\end{table}

Table~\ref{tab:envelope} is not the picture of a cache that is merely too small.
It is the picture of a cache buying the wrong quantity. The direct evidence for
that reading comes from a measurement the design already carries.

\paragraph{The cache reports its own failure rate.}
With the leaf level remote procedure path enabled, a call is dispatched exactly
when the cache could not resolve the leaf. Its rate therefore counts leaf level
cache failures per operation, and it requires no inference from throughput.
Figure~\ref{fig:dexrpc} plots it.

\begin{figure}[t]
  \centering
  \includegraphics[width=0.94\linewidth]{fig_dex_rpc.pdf}
  \caption{Leaf level cache failures per operation across the cache sweep. Flat
  means more memory bought nothing. Both uniform workloads are flat across an
  eightfold increase, the scan row to four decimal places. Both Zipfian workloads
  decline, because there a hot set exists for the cache to hold.}
  \label{fig:dexrpc}
\end{figure}

Under uniform access every lookup misses its leaf at 64, 128 and 256 megabytes,
and 99.7\% still do at 512. Uniform scans are flatter still: the rate holds at
1.3748 to four decimal places across the entire sweep. Under Zipfian access the
same measurement declines as it should, from 0.82 to 0.54 on lookups, because
skew produces a hot set and the lever can capture it. The uniform rows are not a
small effect that more budget would grow. They are a lever reporting that it is
disconnected from the outcome.

The arithmetic behind this is simple. Fifty million keys produce millions of
leaves, and under uniform access every one is equally likely, so there is no hot
core for budget to convert into hit rate. DEX's own admission policy, which
refuses nine leaves in ten, is the design reaching the same conclusion
empirically and responding by tuning the failing lever rather than reaching for
the other one.

\paragraph{The same story in crossings.}
Figure~\ref{fig:dexcrossings} puts it in the unit the fabric charges for.
Caching alone takes uniform scans from 22.3 crossings per operation to 16.9, a
32\% reduction for eight times the memory.

\begin{figure*}[t]
  \centering
  \includegraphics[width=\textwidth]{fig_dex_crossings.pdf}
  \caption{Network crossings per operation, counting one sided reads and remote
  procedure calls together. Note the independent vertical scales: within a panel
  the quantity to read is the gap between the curves, and across panels it is how
  much larger the scan remainders are. Caching alone removes crossings slowly and,
  on uniform scans, barely at all. The second lever removes a roughly constant
  fraction, two to three times, at every cache size.}
  \label{fig:dexcrossings}
\end{figure*}

\subsection{Making the venue a runtime decision}
\label{sec:dex:add}

The change is small and it converts one event. In the design as built a leaf
level miss pulls a page across the network, installs it, and reads eight bytes
out of it. We make that same miss dispatch a request that the memory node answers.
The descent above the leaf, the cache, the admission policy and the binary are
otherwise unchanged.

For a point lookup the memory node searches the leaf in place and returns the
value, so one request replaces one page transfer and the page is never
materialised on the compute node.

For a range scan the change is larger in kind, and it is the contribution rather
than a configuration. Instead of asking for one leaf, the compute node sends the
number of keys still outstanding, and the memory node walks forward along the
sibling chain in its own local memory, gathers matching pairs from as many leaves
as the request needs, and packs them into a contiguous region that the compute
node retrieves with a single read. A scan that spanned a dozen leaves cost a
dozen serial crossings and now costs a constant number, independent of how many
leaves it covers. If the range extends past the last key returned, the traversal
resumes and the remainder is served the same way. This is the mechanism the
decomposition into repeated lookups precludes, and it is available to any index
that can walk its own leaves.

\paragraph{The gate is the whole idea.}
The memory node is engaged only on a miss, so its CPU demand scales with the
\emph{miss} rate rather than with the operation rate. Read as knobs, cache and
memory side execution are substitutes, since more cache means the remote path
fires less often. Read as one structure they are complements: more cache does not
idle the remote path, it \emph{shrinks the bill} that path sends. That is the
difference between this composition and an unconditional pushdown design, which
would pay one request per operation forever regardless of what the cache already
held.

The structure is hybrid along two axes at once. Across the tree it is one sided
above the leaves and result shipping at the leaves, so the inner path remains a
pure caching workload with an idle memory node and only the one level that
caching handles worst is offloaded. Within a single scan it is hybrid per
operation: the resident prefix of a range is scanned locally and costs the memory
node nothing, and only the uncached tail is offloaded.

\subsection{What the second lever buys}
\label{sec:dex:results}

\begin{figure*}[t]
  \centering
  \includegraphics[width=\textwidth]{fig_dex_self.pdf}
  \caption{DEX against itself across the cache sweep. Blue is the design exactly
  as built; green makes the execution venue a runtime decision and differs by one
  setting on the same binary. Top row throughput, higher is better; bottom row
  99th percentile latency, lower is better. The blue curves are what eight times
  the memory buys; the vertical gap at any point is what the second lever buys at
  the same memory.}
  \label{fig:dexself}
\end{figure*}

Figure~\ref{fig:dexself} is the result. At the largest cache the gain is
$1.95\times$ and $1.71\times$ on lookups and $2.62\times$ and $2.50\times$ on
scans, and every one of the sixteen cells improves its tail, by 31\% to 59\%.
Both axes move favourably everywhere, so this is a relocation rather than a
trade.

\paragraph{Against the design's own ceiling.}
The strict test is to take the best throughput the unmodified design reaches at
\emph{any} cache size and count how many cells the second lever puts above it.
Thirteen of sixteen clear that bar on throughput, fourteen on tail, and
\textbf{thirteen on both simultaneously}. All eight scan cells clear it. The three
that do not are lookups at the smallest budgets, where the caching lever still has
room to work.

\paragraph{Cache equivalence.}
Figure~\ref{fig:dexequiv} states the provisioning consequence. On scans the gated
configuration at 64 megabytes beats caching alone at 512 on both axes at once, so
an eighth of the memory buys a strictly better point and the equivalent cache for
the caching lever lies beyond the end of the sweep. On uniform lookups the
crossing is inside the sweep: a quarter of the memory suffices. Only on Zipfian
lookups does the caching lever hold its own, which is the expected place for it to
do so, because skew is what makes a working set cacheable at all.

\begin{figure}[t]
  \centering
  \includegraphics[width=0.94\linewidth]{fig_dex_equivalence.pdf}
  \caption{Cache equivalence. Each bar is the throughput of both levers at 64\,MB
  divided by that of caching alone at 512\,MB, so a bar past the dashed line means
  an eighth of the memory bought a better result. Both scan workloads clear it.}
  \label{fig:dexequiv}
\end{figure}

\paragraph{Why the effect is this large.}
A miss costs 14.25 microseconds when it is a page and 7.27 when it is an answer,
and the 99th percentile of an offloaded operation is 14.0 microseconds, roughly
half the 26.5 microsecond floor the design reaches anywhere in the sweep by
rationing. The mechanism the design declined to build is, on its own, faster than
anything rationing achieved. Cache changes how often the toll is paid; it cannot
change the toll.

\subsection{What it costs}
\label{sec:dex:cost}

The bill is four memory side service threads against thirty six compute threads,
and two properties make it defensible.

First, it is gated by misses rather than by operations, so it shrinks as the
compute side becomes self sufficient. More cache does not strand the mechanism;
it reduces what the mechanism is asked to do.

Second, it saturates early. With the remote path disabled the memory node is
idle and throughput is flat in thread count, which is the control confirming that
the threads do nothing without the path. With it enabled, two service threads
under serve the request stream and four do not: lookup tails fall from 22.5 to
13.5 microseconds between two and four threads and then stop moving at six. Scan
tails do not respond to thread count at all, staying within one histogram bucket
across two, four and six threads, which places the scan tail in the cost of the
offloaded traversal itself rather than in queueing for a service thread.

\subsection{What DEX should have done}
\label{sec:dex:should}

Three critiques follow from the measurements, and each is a design decision
rather than an implementation defect.

\paragraph{The venue decision should have been regime aware.}
DEX's cost model fires at a node miss and compares one request against one
descent. It never asks whether the caching lever is succeeding for an operation
class at all, so when leaf caching degrades the design keeps making locally
reasonable per miss decisions while the aggregate regime has changed underneath
it. The signal it would need is not subtle and the design already produces it:
a leaf failure rate that does not move as the cache grows
(Section~\ref{sec:dex:sat}).

\paragraph{The admission constant should have been read as a diagnosis.}
Refusing nine leaves in ten is the correct empirical response to leaves not being
worth caching. But it is a symptom, and the same observation argues that those
accesses should be served by the other lever, which the design already possessed.
The two findings sit in adjacent parts of the design and are never connected.

\paragraph{The scan path should not have been traded for pointer simplicity.}
Dropping leaf links keeps pointer management uniform, at the price of turning
every multi leaf scan into a sequence of independent descents. That price falls
on the one operation class where the caching lever is structurally unavailable,
because a scan streams through leaves it will never revisit. A leaf run mechanism
used only on the offload path, never in the cached structure, keeps the
simplification and still allows one reply to carry a run.

% ===========================================================================
\section{CHIME}
\label{sec:chime}

\subsection{The design as built}

CHIME is the harder subject and the better control, because it economises in a
different direction and because its whole structure is built around the memory
node doing nothing.

Its compute cache holds internal nodes only. This is deliberate: excluding leaves
avoids cross node coherence entirely, at the price that the leaf is always a
remote access even when the cache resolved the whole descent. Everything else in
the design makes that unavoidable access smaller. A hopscotch layout confines a
key to a short neighbourhood so a reader fetches a segment rather than a whole
leaf. A hotspot structure remembers which slot a hot key occupies so the read
shrinks to a single entry. Replicated per group metadata makes a partial read self
describing. Per cacheline version bytes let a one sided reader detect a torn read
without a lock.

Two further choices set how badly the design behaves when its bet fails. The
internal fanout is small, chosen so that an internal node is smaller than a leaf,
which makes the tree deep: in our runs the root sits seven levels above the
leaves. And the internal index is roughly 90 to 100 megabytes at this geometry.
Together they mean a cache hit costs one leaf access while a cache miss costs
about seven crossings, so the design is excellent above its residency boundary and
falls a long way below it.

\subsection{Adding both levers}
\label{sec:chime:add}

We add memory side execution on the same gate as in DEX, and we add a compute
side cache for leaves so that the caching lever can reach the last crossing as
well. Both are runtime settings, so one build serves every cell of the sweep.

\paragraph{Encoding agnostic execution.}
CHIME's leaves are not plain records. They are version encoded per cacheline and
their metadata is scattered across entry groups, precisely so that a one sided
reader can validate them without a lock. A memory node that wants to answer a
request about such a leaf cannot simply point at it. It must take a private
snapshot, check that the snapshot is internally consistent, undo the scattering,
and retry if a concurrent writer intervened. The snapshot has to come first:
comparing version bytes in place would let a writer land another cacheline between
the check and the read.

This is the general point. Memory side execution does not require an index to
store plain records; it requires the memory node to speak the index's own
encoding. The traversal logic above that decoding step is identical in both of our
subjects. Only the decoding differs, and it is confined to a single function. An
index whose node format was designed to make memory side CPU unnecessary can still
use memory side CPU.

\paragraph{Structural gating.}
CHIME caches only the level immediately above the leaves, which makes the miss
gate a property of the structure rather than a policy: a cache hit resolves the
descent completely and therefore cannot reach the offload path at all. The gate
is validated independently by comparing the observed offload fraction against the
miss rate predicted from cache occupancy, and the two agree at three separate
budgets, at 63\% against a predicted 60\%, 65\% against 62\%, and 83\% against
81\%. The memory node is busy because the cache is missing, not because the policy
is loose.

\paragraph{Leaf caching and its coherence.}
A leaf served out of local memory is the one read in CHIME that is not self
validating, so it needs a coherence protocol. Each leaf carries a small,
globally unique, never reused stamp that a writer publishes under the leaf's own
lock before it modifies any data. A reader may serve a cached copy only while a
single small read of the leaf's lock and stamp still shows the leaf unlocked and
carrying the stamp the copy was taken at. Because the writer's operations travel
in order on one connection, this makes the check conclusive rather than
heuristic, and because a filling read can close the same check around the data it
fetches, populating the cache costs the round trip the uncached read would have
cost anyway. The cache is never weaker than the one sided path it replaces.

\subsection{The lever map}
\label{sec:chime:map}

The sweep crosses the residency boundary, which is what Arm one could not do, and
the analysis step that makes it legible is to index results by the \emph{inner}
cache rather than by the total. Inner cache is what the caching lever spends on
the descent, and the leaf arm halves it. Re indexed that way the two arms become
commensurable: a leaf enabled cell at 128 megabytes total and a leaf disabled cell
at 64 megabytes total both run a 64 megabyte inner cache, and they measure 0.645
and 0.632 million operations per second, within two percent of each other despite
one holding twice the total memory. Performance here is a function of the inner
budget and the total is close to irrelevant.

\begin{figure*}[t]
  \centering
  \includegraphics[width=\textwidth]{fig_lever_map.pdf}
  \caption{The lever map. Compute node throughput against inner node cache, which
  is what the caching lever spends. The shaded band is where the internal index
  stops fitting. Caching raises the right hand plateau and does almost nothing to
  the left of the band; memory side execution raises the left hand floor and is
  inert to the right of it. The two levers act on disjoint parts of the curve.}
  \label{fig:levermap}
\end{figure*}

Figure~\ref{fig:levermap} is the result, and each lever has a signature.

\textbf{Caching raises the plateau.} Adding the leaf cache lifts Zipfian point
throughput from 2.857 to 3.999 million operations per second, a 40\% gain, and
batched Zipfian scans from 0.500 to 0.594, a 19\% gain. To the left of the band it
does essentially nothing: 1.091 becomes 1.250.

\textbf{Memory side execution removes the cliff.} At a 64 megabyte inner cache,
Zipfian point throughput goes from 1.091 to 2.727. To the right of the band it is
\emph{exactly} inert, holding 2.857 across three consecutive cache points. That
flatness is the gate working as specified.

\textbf{Together they exceed what either reaches.} At a 32 megabyte inner cache
with 32 megabytes of leaf cache, Zipfian point throughput reaches 3.157. The
unmodified design's ceiling is 2.857 and it needs at least a 128 megabyte inner
cache to hold it. The pair therefore delivers 1.11 times the best the design ever
reaches, on half the total budget and a quarter of the inner budget.

\paragraph{One lever can make things worse, and the other bounds it.}
At a 128 megabyte total budget the leaf arm runs a 64 megabyte inner cache, below
the band, and uniform point throughput falls from 2.307 to 0.645. That is a
factor of 3.6 lost by spending budget on the wrong lever. With memory side
execution enabled the same cell recovers to 1.714. The second lever is therefore
not merely additive: it bounds the downside of using the first one aggressively,
which is what makes an aggressive caching policy safe to adopt at all.

\paragraph{Where it does not hold.}
The headline is a Zipfian result. Under uniform access both levers at a 32
megabyte inner cache reach 1.666 against an unmodified plateau of 2.307, so the
pair removes the cliff but does not exceed the ceiling. And on uniform scans leaf
caching is a net loss at every budget we swept, costing 21\% at the largest,
because a uniform scan streams through leaves it will never revisit and there is
nothing for that lever to hold.

\subsection{A falsification test}
\label{sec:chime:falsify}

The account underlying all of this is that crossings and not bytes are the
currency. That claim forbids something specific: an optimisation that shrinks the
bytes a crossing carries, applied where the baseline already spends exactly one
crossing, must not help, and must not be rescued by a high hit rate either.

CHIME's scan path is exactly that setting, since it spends one read per covered
leaf. Our first leaf cache did shrink those reads, from roughly a kilobyte to a
small validation probe on a hit, a factor of about forty, while leaving the
crossing count untouched: the probe is still issued per covered leaf. The account
predicts a regression, and the measurement delivered one of 41\% on uniform scans.
More tellingly it also lost 15\% on Zipfian scans \emph{at a 59\% hit rate}, which
rules out the obvious alternative explanation that the cache was simply too small.

\begin{figure}[t]
  \centering
  \includegraphics[width=0.94\linewidth]{fig_scan_repair.pdf}
  \caption{The falsification test and its repair. Shrinking the bytes per covered
  leaf by roughly forty times while leaving the crossing count alone loses
  throughput, and loses it even where the hit rate is 59\%. Removing crossings
  instead turns the Zipfian arm from a regression into a gain.}
  \label{fig:scanrepair}
\end{figure}

The repair follows from the same account. The set of leaves a scan covers is
known before any network operation is issued, so the validations need not be
serial: one batch validates every resident leaf and a second fetches the rest,
taking a twelve leaf scan from twelve crossings to two. Figure~\ref{fig:scanrepair}
shows the result. The Zipfian arm moves from a 15\% loss to a 19\% gain and the
uniform arm from a 41\% loss to a 21\% loss. The leaf arm gains between 35\% and
40\% in all four cells, and the near uniformity of that gain is itself evidence: a
crossing count fix should appear as a flat multiplier that does not depend on hit
rate, and it does.

The point path is untouched by the repair and reproduces across two independent
six hour sweeps to four decimal places, so run to run noise where the index fits
is negligible and the scan movements are signal. In the stressed cells noise runs
five to twelve percent, and small differences there should not be read.

\subsection{What CHIME should have done}

\paragraph{It should not have bet the design on the index fitting.}
Inner node caching gives a plateau that is flat and good and a cliff that is a
factor of 3.7 deep, with no second lever to catch it. The data shows the catch is
cheap: memory side execution restores most of the loss and is exactly inert above
the boundary, so it costs nothing when it is not needed. A system with one lever
must be provisioned for its worst case; a system with two can be provisioned for
its common case.

\paragraph{It should have priced the exclusion of leaves rather than assuming it.}
Excluding leaves avoids coherence entirely, and the price is one guaranteed
crossing per operation forever. Our stamp protocol shows that price was
mis estimated, since coherence for a cached leaf costs one small probe folded into
a crossing the read needed anyway. Stated carefully, because our own map
complicates it: the leaf lever pays only for point lookups under skew and for
batched Zipfian scans, and it is a loss on uniform scans. The recommendation is
not that leaves should be cached. It is that the coherence argument alone never
settled the question, and the answer is regime dependent.

\paragraph{It should have batched its scan path regardless of any lever.}
The set of covered leaves is known before any network operation is issued, yet the
unmodified path issues one read per leaf, serially. The information needed to
batch was already in hand and unused. This is the cheapest structural improvement
available in the design, it requires no memory side CPU and no coherence protocol,
and it is what made leaf caching viable for scans at all.

\paragraph{It should have reported against its residency boundary.}
Throughput against total cache hides a factor of 3.7 cliff and makes the two
caching arms incomparable; throughput against inner cache makes the structure
legible immediately. The distinction between nominal and effective index cache
matters for the same reason, since a hotspot structure silently claims part of the
nominal budget.

% ===========================================================================
\section{Discussion}
\label{sec:discussion}

\subsection{Three design rules}

The two studies converge on the same three rules, and they are properties of
execution venue rather than of B$^+$ trees.

\paragraph{Ship answers, not pages, but only where caching loses.}
The inner path is where caching wins and should stay one sided and memory node
free. The leaf level is where caching structurally cannot win, because leaf
residency scales with the data rather than with the index. Choosing the venue per
level is what keeps the memory node idle for the majority of the work.

\paragraph{Gate the remote venue on evidence the index already has.}
A miss is a free and exact signal that the compute side cannot answer locally.
Conditioning remote execution on it makes the cost self limiting, which is what
separates this from unconditional pushdown: cache and memory side execution become
complements, and adding cache reduces the memory node bill instead of stranding
the mechanism.

\paragraph{Make the venue decision per operation, not per system.}
The within scan behaviour is the sharpest evidence. A scan whose prefix is
resident is served entirely locally and costs the memory node nothing; only the
uncached tail is offloaded.

\subsection{When the second lever is the wrong answer}

It is inert where caching succeeds, and that is by construction rather than by
luck: in CHIME above the residency boundary it holds one number across three
consecutive cache points. Where a hot set exists and fits, rationing is the right
lever, and our own measurements show the cached configuration winning outright on
Zipfian lookups. A design that always offloads would spend memory node CPU in
exactly the regime where it buys nothing.

\subsection{Limitations}
\label{sec:limits}

\paragraph{Scan tails.}
The second lever improves scan tails substantially but they remain several times
lookup tails, because an offloaded scan walks several leaves on the memory node
and that work sets the floor. Separately, CHIME's leaf caching arm gains
throughput on Zipfian scans while its tail moves the wrong way, plausibly because
batching replaces several small stalls with one larger one. We report this rather
than quoting the throughput alone.

\paragraph{Reads only.}
Both sweeps are read only, so nothing here bounds insert, update or delete
behaviour, and the leaf coherence protocol is exercised but idle: no cached copy
was ever refused during a measured run. A write mixed workload is the measurement
that would give that protocol meaning.

\paragraph{Scale out.}
Both studies use one compute node and one memory node. Self limiting per compute
node is not the same property as bounded across many, and the memory tier is
fixed while compute nodes are not. The gate predicts that aggregate demand should
scale with total misses rather than with total operations, which is a smaller and
more compressible quantity, and that prediction has an obvious falsifier.

\paragraph{No cross system comparison.}
For the reasons in Section~\ref{sec:method}. Adding one requires the four
equalisations stated there to be fixed and reported, the same latency statistic on
both sides, and a sweep run for that purpose rather than assembled from sweeps
that happened to exist.

% ===========================================================================
\section{Related Work}
\label{sec:related}

Disaggregated ordered indexes have converged on compute side caching with one
sided access. Sherman caches only the lowest inner levels and accepts one remote
access per leaf; SMART makes a similar trade for a trie. Both exclude leaves from
the cache for the coherence reason discussed in Section~\ref{sec:bg}, and neither
executes index logic on the memory node. DEX extends caching to leaves under a
path aware policy with probabilistic admission, and adds a cost model that
offloads point operations at the cache boundary, which makes it the closest prior
work and the reason we treat it as the primary subject rather than as a foil.
CHIME attacks trip shape instead of trip frequency. Designs that push computation
into the memory tier exist outside the index literature, generally with
unconditional policies, which is the behaviour our gate is contrasted against.
What is absent from all of them is a scan served as a run rather than as a
sequence of independent descents, and a venue decision conditioned on evidence
that the cache has already failed.

% ===========================================================================
\section{Conclusion}
\label{sec:conclusion}

Taken exactly as built, both of our subjects walk along a curve that more memory
moves them only a little way down, and in neither case is the cause a cache that
is too small. It is that a cache changes how often the remote toll is paid and
never what the toll is. DEX says so itself: at the largest cache only 0.31\% of
uniform lookups are served without leaving the compute node, and the rate at which
its leaf path fires does not fall at all under uniform access across an eightfold
sweep. CHIME says so differently, by falling a factor of 3.7 in one step when its
internal index stops fitting.

Converting one event, a leaf level miss, from a page transfer into an answer, and
gating that conversion on the miss itself, puts thirteen of sixteen DEX cells past
the best that design reaches anywhere in its sweep, on throughput and tail at
once, and on scans buys a better operating point at one eighth of the memory. In
CHIME it separates the two levers by regime, and their combination exceeds the
unmodified ceiling on a quarter of the inner cache needed to reach it.

The account behind those results is falsifiable and we tried to falsify it. An
optimisation that removed bytes without removing crossings lost 41\% of
throughput, and lost even where more than half of its lookups hit. The same
mechanism, given a path that removes crossings instead, gains 19\%. Crossings are
the currency, the execution venue is what sets them, and it should be a runtime
decision rather than a constant chosen when the index was designed.

\balance
\end{document}
